% ============================================================================
% Online Appendix (Supplementary Material): a self-contained PDF submitted
% separately from the main manuscript. It collects the external empirical
% anchoring of the social-instability parameters (Section C) and the extended
% robustness runs (Section D). Every number is produced by the simulation and
% analysis scripts in the repository.
% ============================================================================
\documentclass[12pt]{article}
\usepackage[margin=1in]{geometry}
\usepackage{amsmath,amssymb,amsthm}
\usepackage{booktabs,array}
\usepackage{graphicx}
\usepackage[utf8]{inputenc}
\usepackage[T1]{fontenc}
\usepackage{lmodern}
\usepackage{xurl}
\usepackage{microtype}
\graphicspath{{figures/}}

\theoremstyle{plain}
\newtheorem{proposition}{Proposition}

\begin{document}

\begin{center}
{\Large\bfseries Online Appendix}\\[4pt]
Anatomy of a Smooth Transition: A Task-Based Macroeconomic Model of AGI
Displacement with Fiscal and Social Safety Nets\\[2pt]
DUAN Fang, Huaqiao University
\end{center}

\tableofcontents
\bigskip

% ============================================================================
\section{External empirical anchoring of the social-instability parameters}
% ============================================================================
The social-instability channel contains three reduced-form parameters: the
social-crash cutoff $\varsigma^{th}$ (baseline $0.50$), the stress coefficient
$\alpha_s$ (baseline $0.60$), and the sustained-exceedance window $J$
(baseline $10$ periods, one model period being one year). We search the
empirical literature for mappable evidence for each, distinguish direct
anchors from indirect directional evidence, and state the mapping assumptions
explicitly. Table~\ref{tab:anchor-status} summarizes the anchoring status.

\begin{table}[h]
\centering\small
\caption{External empirical anchoring status of the three social-instability parameters.}
\label{tab:anchor-status}
\begin{tabular}{@{}
>{\raggedright\arraybackslash}p{0.8in}
>{\raggedright\arraybackslash}p{0.7in}
>{\raggedright\arraybackslash}p{1.7in}
>{\raggedright\arraybackslash}p{2.6in}@{}}
\toprule
Parameter & Baseline & Anchoring status & Nature of evidence \\
\midrule
$\varsigma^{th}$ & 0.50 & no direct external anchor & directional/continuous probability only; no normalized cutoff \\
$\alpha_s$ & 0.60 & partially anchored & illustrative reverse-engineered range $\approx1.0$--$6.6$; baseline 0.60 is below this range, a conservative setting \\
$J$ & 10 & partially anchored & mean/upper-tail scale (crisis 8--10 yrs; civil-war mean 11.1 yrs) \\
\bottomrule
\end{tabular}
\end{table}

\paragraph{C.1 The cutoff $\varsigma^{th}=0.50$: no direct external anchor.}
No study provides a population-dependency threshold that maps onto a $0.50$
cutoff of the $[0,1]$ normalized stress variable. The authoritative literature
consistently reports continuous probability or dose--response relationships
rather than structural thresholds replicable across samples. Using rainfall as
an instrument across 41 African countries over 1981--1999, Miguel, Satyanath
and Sergenti~\cite{MSS2004} identify that a five-percentage-point drop in real
GDP growth raises the probability of conflict in the following year (at least
25 battle deaths) by about 12 percentage points, a relative increase of
roughly 50\%. Ponticelli and Voth~\cite{PV2020}, across 25 European countries
over 1919--2008, find that a one-standard-deviation spending cut (about 2.8
percentage points of GDP) adds roughly 0.4 unrest events per year, with the
event count nearly doubling as cuts widen from 1\% to 10\% of GDP; only a
descriptive soft cutoff around a 2\%-of-GDP cut appears, not an identified
structural threshold. Using natural disasters as instruments,
Bellemare~\cite{Bellemare2015} identifies that a one-unit rise in the global
food price index adds about 0.99 food-related unrest reports per month, while
price volatility has no significant effect. Earlier work by Collier and
Hoeffler~\cite{CH2004}, Fearon and Laitin~\cite{FL2003}, and Alesina and
Perotti~\cite{AP1996} likewise reports continuous marginal effects. The only
explicit numerical cutoff, due to Lagi, Bertrand and Bar-Yam (2011,
arXiv:1108.2455, not peer reviewed), clusters protests when the FAO food price
index exceeds 210; but that cutoff is a price-index level rather than a
population dependency ratio, and mapping it onto $\varsigma^{th}=0.50$ would
require an over-long chain of normalization assumptions, so we do not adopt it
as a direct anchor.

\paragraph{C.2 The stress coefficient $\alpha_s=0.60$: partially anchored, with an identifiable range.}
The sign with which unemployment or relief dependence raises property crime
and social unrest is robust across U.S., European and cross-country panels,
ruling out $\alpha_s\leq0$. The literature reports crime semi-elasticities to
unemployment of roughly $+1.3\%$ to $+4\%$ per percentage point of
unemployment, concentrated in property crime with weak effects on violent
crime. Raphael and Winter-Ebmer~\cite{RWE2001} find that a one-percentage-point
fall in unemployment lowers property crime by about 1.6\%--2.4\%;
Lin~\cite{Lin2008} identifies a 1.8\% OLS and 4.0\% 2SLS rise in property
crime per extra percentage point of unemployment; Gould, Weinberg and
Mustard~\cite{GWM2002} find that the 23.3\% decline in low-skill male wages
over 1979--1993 accounts for 43\% of the rise in property crime; Buonanno,
Drago and Galbiati~\cite{BDG2014}, across about 20 developed countries over
1979--2010, identify a roughly 1.3\% rise in total crime and 1.7\% rise in
homicides per extra percentage point of unemployment; and Foug\`ere, Kramarz
and Pouget~\cite{FKP2009} identify a significant positive effect of youth
unemployment on property crime in France; see also the surveys by
Freeman~\cite{Freeman1999} and Draca and Machin~\cite{DM2015}. Mapping these
semi-elasticities $\beta_{crime}\approx\partial\log C/\partial u$ onto
$\alpha_s=\partial\,\text{stress}/\partial r$ requires three layers of
assumptions, and we keep every quantity in consistent units: (A) the normalized
range linking the stress index to disorder indicators, setting
$\text{stress}=1$ to correspond to a 3--10-fold amplification of crime over
baseline (mapping slopes $f'$ of about 0.43--0.91); (B) the conversion factor
$k=r/u$ between the relief share $r$ and unemployment $u$; and (C) a local
linearization treating the semi-elasticity as approximately constant over the
working range. On (B), rather than import an external value, we use the model's
own accounting: relief recipients are the unemployed not absorbed by public
employment, $r=(1-p_E)u$, so at the baseline $p_E=0.45$ the conversion is
$k=1-p_E=0.55$ (verified directly from the simulated paths). On units, when
$u$ is measured on the $[0,1]$ scale the semi-elasticity entering the chain is
$\beta_{crime}=1.3$--$4$ (the per-percentage-point $1.3\%$--$4\%$ figures
restated per unit of the fraction $u$). Substituting into
$\alpha_s\approx f'\beta_{crime}(1/k)$ with $k=0.55$ gives an illustrative
reverse-engineered interval of roughly $1.0$--$6.6$ (the lower value
$0.43\times1.3/0.55\approx1.02$, the upper $0.91\times4/0.55\approx6.62$).
We read this as a wide, normalization-dependent illustration rather than an
empirical estimate: the interval spans more than a factor of six and moves
almost one-for-one with the assumed stress-to-disorder mapping. The baseline
$0.60$ lies below even the lower end of this interval, so the calibration
deliberately errs on the conservative side, assigning a smaller stress
response than the reverse engineering would tolerate; it is a scenario setting
rather than a value the interval identifies.

\paragraph{C.3 The sustained-exceedance window $J=10$: partially anchored at mean/upper-tail scale.}
No study directly estimates the window of consecutive stress exceedance that
triggers a crash. Annual-scale evidence brackets ten years along two
independent lines. First, after systemic banking crises, per-capita GDP returns
to its pre-crisis peak in about eight years on average (median 6.5 years), with
45\% of episodes experiencing a second dip~\cite{RR2014}; Cerra and
Saxena~\cite{CS2008} find highly persistent, near-permanent output losses
after banking crises, with later estimates placing the average post-crisis GDP
loss at about ten years. Second, civil wars last 11.1 years on average (median
7.1 years)~\cite{Fearon2004}, and the hazard estimates of Collier, Hoeffler
and S\"oderbom~\cite{CHS2004} average just over seven years, with wars begun
in the 1980s lasting 105 months (about 8.75 years). On the other side, protest
waves and crime shocks persist only months to 2--5 years, markedly shorter
than ten years. $J=10$ is therefore a defensible mean/upper-tail anchor
aligned with chronic post-crisis slumps and average civil wars, but it lies
above the median durations and should be read as such rather than as a median
anchor.

\paragraph{C.4 Parameter decision and empirical-range robustness.}
Because none of the three parameters has a precise, directly identified
empirical replacement, we keep the baseline values but state their standing
accurately: assigning a value to $\varsigma^{th}$ would lack mappable threshold
evidence; $\alpha_s=0.60$ is a conservative setting that sits below the
$\approx1.0$--$6.6$ reverse-engineered interval rather than inside it; and
$J=10$ is a supported mean/upper-tail value. An empirical-range sweep extends
the $\alpha_s$ scan across the full $0.30$--$6.62$ interval alongside a
directional $\varsigma^{th}$ scan. The analytic boundaries are
$\alpha_s^{crit}=0.50/((1-0.45)u^{max})=0.909$ (with $u^{max}=1$) and
$\varsigma^{th}_{crit}=0.33$, and the numerical sweep locates the first social
crash at $\alpha_s=0.91$ (smooth at $0.90$) and between
$\varsigma^{th}=0.30$ and $0.33$ (smooth at $0.33$); $\min C$ stays $0.357$
throughout because the social channel does not feed back into the economic
paths. Two results bound the channel. First, the ranking is robust across the
entire interval: the pure transfer crashes at every $\alpha_s$, whereas the
corner stays smooth through $\alpha_s=0.90$. Second, the corner's own
smoothness is not unconditional: under the upper portion of the
reverse-engineered interval ($\alpha_s\geq0.91$), the corner also triggers the
social channel. The sweep therefore reports both the conservative margin at
the baseline and the parameter region in which even the workfare corner would
fail, rather than presenting the reduced-form channel as settled.

% ============================================================================
\section{Extended robustness runs}
% ============================================================================
This section collects the three extensions summarized in the body: the
functional-form comparison of the share law, the capability-sensitivity sweep,
and the counterfactual that removes the logistic erosion term. All runs are
reproduced by the named analysis scripts.

\paragraph{D.1 Functional form of the share law.}
Equation for the endogenous rate is a modelling choice, so we re-solve the
corner with the probit share, $\omega_t=\Phi(\ell_0+\Gamma_t-\beta_A g_t)$
where $\Phi$ is the standard normal cumulative distribution, which shares the
logistic form's endpoints, monotonicity and smoothness but a different tail
shape. Both forms keep the corner smooth and without a real crossing. Measured
over the policy-driven horizon $t\geq1$, the effective capture peaks at
$0.121$ under the logistic and $0.041$ under the probit. The buffer conclusion
is nevertheless unchanged, with $\min_t C_t$ moving only from $0.3562$ to
$0.3571$, a gap below $0.001$. The binding constraint remains skill
depreciation rather than the choice of link function, so the functional form
is not the margin that decides the outcome. This is reproduced by
\texttt{analysis/check\_link\_function.py}.

\paragraph{D.2 Capability-sensitivity sweep.}
On the $T=200$ horizon, varying $\beta_A\in\{0,1,3\}$ separates a
mid-horizon crossing from a durable reinstatement; the critical value is
$0.5667$. At $\beta_A=0$, with no capability erosion, the strong floor crosses
at $t=29$ and the effective capture peaks at $0.731$ and stays above the
threshold, so the crossing is durable in the sense that it is not reversed;
the market employment share ends the horizon at a positive $0.344$ without
sliding back, although it still declines at a small negative rate rather than
settling. At $\beta_A=1$ the strong floor crosses at $t=35$, with a peak of
$0.640$ above the critical value, but continued erosion pulls the effective
capture back below the threshold later, so the employment share keeps declining
to $0.172$: a mid-horizon crossing that defends consumption and employment
over the short run without a durable reinstatement. At $\beta_A=3$ the peak
falls to $0.486$, below the critical value, there is no crossing and the
employment share reaches zero. The corner remains smooth throughout and its
minimum consumption stays in $0.356$--$0.357$, while the strong floor's
minimum consumption stays in $0.433$--$0.456$ at every reported $\beta_A$, so
the fiscal buffer defends consumption at every setting even though a durable
restoration of the labor engine occurs only in the no-erosion limit. The
baseline $\beta_A=1$ is the conservative middle setting, not an estimated
value.

\paragraph{D.3 Counterfactual: the collapse is not built into the logistic term.}
The law contains a monotone term $-\beta_A g_t$ with $g_t=\log(\mathrm{cap}_t
/A_0)$ accumulating over the horizon, so a reader could ask whether the
labor-share collapse is imposed by the functional form rather than generated
by the model. The direct counterfactual sets $\beta_A=0$, which removes the
erosion channel and leaves only the policy stock and the task accounting, and
re-runs the same engine with every other parameter fixed. The no-policy path
still crashes, with $\min C=0.0575$ and the employment share absorbed to
zero, and the financed corner still stays smooth with $\min C=0.356$. The
collapse therefore follows from the task-accounting condition
$\omega\sigma\,dN<\kappa_I A$ of Theorem~1 in the main manuscript, which holds
even when the reinstatement rate is not being eroded, rather than from the
monotone logistic term. The $\beta_A$ coefficient instead governs the
durability of a strong-policy reinstatement, the terminal share $0.344$ versus
$0.172$ reported above, not whether the no-policy path crashes. This
counterfactual is reproduced by
\texttt{analysis/check\_capability\_counterfactual.py}.

% ============================================================================
\begin{thebibliography}{99}
% ============================================================================

\bibitem{MSS2004} E. Miguel, S. Satyanath and E. Sergenti, Economic Shocks and Civil Conflict: An Instrumental Variables Approach, \emph{Journal of Political Economy} 112(4), 725--753, 2004, DOI: \texttt{10.1086/421174}.

\bibitem{PV2020} J. Ponticelli and H.-J. Voth, Austerity and Anarchy: Budget Cuts and Social Unrest in Europe, 1919--2008, \emph{Journal of Comparative Economics} 48(1), 1--19, 2020, DOI: \texttt{10.1016/j.jce.2019.09.007}.

\bibitem{Bellemare2015} M. F. Bellemare, Rising Food Prices, Food Price Volatility, and Social Unrest, \emph{American Journal of Agricultural Economics} 97(1), 1--21, 2015, DOI: \texttt{10.1093/ajae/aau038}.

\bibitem{RWE2001} S. Raphael and R. Winter-Ebmer, Identifying the Effect of Unemployment on Crime, \emph{Journal of Law and Economics} 44(1), 259--283, 2001, DOI: \texttt{10.1086/320275}.

\bibitem{Lin2008} M.-J. Lin, Does Unemployment Increase Crime? Evidence from U.S. Data 1974--2000, \emph{Journal of Human Resources} 43(2), 413--436, 2008, DOI: \texttt{10.3368/jhr.43.2.413}.

\bibitem{GWM2002} E. D. Gould, B. A. Weinberg and D. B. Mustard, Crime Rates and Local Labor Market Opportunities in the United States: 1979--1997, \emph{Review of Economics and Statistics} 84(1), 45--61, 2002, DOI: \texttt{10.1162/003465302317331919}.

\bibitem{BDG2014} P. Buonanno, F. Drago and R. Galbiati, Response of Crime to Unemployment, \emph{Journal of Contemporary Criminal Justice} 30(1), 29--40, 2014, DOI: \texttt{10.1177/1043986213509023}.

\bibitem{FKP2009} D. Foug\`ere, F. Kramarz and J. Pouget, Youth Unemployment and Crime in France, \emph{Journal of the European Economic Association} 7(5), 909--938, 2009, DOI: \texttt{10.1162/JEEA.2009.7.5.909}.

\bibitem{RR2014} C. M. Reinhart and K. S. Rogoff, Recovery from Financial Crises: Evidence from 100 Episodes, \emph{American Economic Review Papers and Proceedings} 104(5), 50--55, 2014 (NBER Working Paper 19823), DOI: \texttt{10.3386/w19823}.

\bibitem{CS2008} V. Cerra and S. C. Saxena, Growth Dynamics: The Myth of Economic Recovery, \emph{American Economic Review} 98(1), 439--457, 2008, DOI: \texttt{10.1257/aer.98.1.439}.

\bibitem{Fearon2004} J. D. Fearon, Why Do Some Civil Wars Last So Much Longer Than Others?, \emph{Journal of Peace Research} 41(3), 275--301, 2004, DOI: \texttt{10.1177/0022343304043770}.

\bibitem{CHS2004} P. Collier, A. Hoeffler and M. S\"oderbom, On the Duration of Civil War, \emph{Journal of Peace Research} 41(3), 253--273, 2004, DOI: \texttt{10.1177/0022343304043769}.

\bibitem{CH2004} P. Collier and A. Hoeffler, Greed and Grievance in Civil War, \emph{Oxford Economic Papers} 56(4), 563--595, 2004, DOI: \texttt{10.1093/oep/gpf064}.

\bibitem{FL2003} J. D. Fearon and D. D. Laitin, Ethnicity, Insurgency, and Civil War, \emph{American Political Science Review} 97(1), 75--90, 2003, DOI: \texttt{10.1017/S0003055403000534}.

\bibitem{AP1996} A. Alesina and R. Perotti, Income Distribution, Political Instability, and Investment, \emph{European Economic Review} 40(6), 1203--1228, 1996, DOI: \texttt{10.1016/0014-2921(95)00030-5}.

\bibitem{Freeman1999} R. B. Freeman, The Economics of Crime, in O. Ashenfelter and D. Card (eds.), \emph{Handbook of Labor Economics}, vol.~3, ch.~52, Elsevier, Amsterdam, 3529--3571, 1999.

\bibitem{DM2015} M. Draca and S. Machin, Crime and Economic Incentives, \emph{Annual Review of Economics} 7, 389--408, 2015, DOI: \texttt{10.1146/annurev-economics-080614-115808}.

\end{thebibliography}

\end{document}
